\documentclass[10pt]{article}
\usepackage[T1]{fontenc}
\usepackage[utf8]{inputenc}
\usepackage[french]{babel}
\usepackage{amsmath,amssymb,mathrsfs,mathpazo}
\usepackage{geometry}
\geometry{a4paper,margin=2cm}
\usepackage{graphicx,array,tabularx,tabularray,colortbl,multirow,multicol,booktabs,enumitem,pifont,ulem,fancybox,diagbox,cancel}
\usepackage{pstricks,pstricks-add,pst-plot,pst-tree,pst-node,pst-text,pst-func,pst-math,pst-bspline,pst-3dplot,pst-eucl,pst-solides3d}
\let\pscalculate\undefined
\usepackage{pst-bezier}
\usepackage{tikz,pgfplots,tkz-tab}
\usetikzlibrary{arrows,arrows.meta,patterns,calc,positioning,trees,decorations.markings,shapes.geometric,shapes.misc,shapes.arrows,intersections,decorations.pathmorphing,decorations.pathreplacing,backgrounds,shadows,babel}
\pgfplotsset{compat=1.18}
\usepgfplotslibrary{fillbetween}
\usepackage[np]{numprint}
\usepackage{listings,makecell,pas-tableur,twoopt}
\let\barre\undefined
\usepackage[pstricks]{tablvar}
\definecolor{darkgreen}{RGB}{0,100,0}
\providecommand{\red}{\color{red}}
\providecommand{\blue}{\color{blue}}
\providecommand{\green}{\color{green}}
\providecommand{\black}{\color{black}}
\providecommand{\white}{\color{white}}
\providecommand{\gray}{\color{gray}}
\providecommand{\cyan}{\color{cyan}}
\providecommand{\magenta}{\color{magenta}}
\providecommand{\yellow}{\color{yellow}}
\providecommand{\violet}{\color{violet}}
\providecommand{\orange}{\color{orange}}
\providecommand{\euro}{\texteuro}
\providecommand{\EUR}{\texteuro}
\providecommand{\ieme}{\textsuperscript{e}}
\providecommand{\degres}{\ensuremath{^\circ}}
\providecommand{\degree}{\ensuremath{^\circ}}
\providecommand{\vv}[1]{\overrightarrow{#1}}
\providecommand{\Coord}[1]{#1}
\providecommand{\decosix}{}
\providecommand{\starredbullet}{\ensuremath{\bullet}}
\pagestyle{empty}
\setlength{\parindent}{0pt}
\usepackage{adjustbox}\DeclareUnicodeCharacter{2236}{:}\DeclareUnicodeCharacter{2190}{\ensuremath{\leftarrow}}
\providecommand{\R}{\mathbb{R}}
\providecommand{\N}{\mathbb{N}}
\providecommand{\Z}{\mathbb{Z}}
\providecommand{\Q}{\mathbb{Q}}
\providecommand{\D}{\mathbb{D}}
\providecommand{\C}{\mathbb{C}}
\providecommand{\e}{{\mathrm{e}}}
\providecommand{\vect}[1]{\overrightarrow{#1}}
\providecommand{\vectt}[1]{\overrightarrow{\mathrm{#1}}}
\providecommand{\vvt}[1]{\overrightarrow{\mathrm{#1}}}
\providecommand{\np}[2][]{{#2~#1}}
\providecommand{\vv}[1]{{\overrightarrow{#1}}}
\renewcommand{\d}{{\mathrm{d}}}
\providecommand{\degres}{{^\circ}}
\providecommand{\degre}{{^\circ}}
\providecommand{\blue}{\color{blue}}
\providecommand{\red}{\color{red}}
\providecommand{\enskip}{\quad}
\providecommand{\llbracket}{\lbrack\!\lbrack}
\providecommand{\rrbracket}{\rbrack\!\rbrack}
\providecommand{\euro}{{\text{€}}}
\providecommand{\no}{{\text{n°}}}
\providecommand{\interval}[2]{{[#1;#2]}}
\providecommand{\intervFO}[2]{{[#1;#2[}}
\providecommand{\dotfill}{\ldots}

\providecommand{\ds}{\displaystyle}
\providecommand{\textding}[1]{#1}
\providecommand{\Oij}{\(\mathrm{O};\vec{\imath},\vec{\jmath}\)}
\providecommand{\Oijk}{\(\mathrm{O};\vec{\imath},\vec{\jmath},\vec{k}\)}
\providecommand{\Ouv}{\(\mathrm{O};\vec{u},\vec{v}\)}
\providecommand{\intervFF}[2]{\ensuremath{[#1;#2]}}
\renewcommand{\euro}{\ifmmode\text{\texteuro}\else\texteuro\fi}
\DeclareUnicodeCharacter{00B0}{\ifmmode^{\circ}\else\textdegree\fi}
\DeclareUnicodeCharacter{20AC}{\ifmmode\text{\texteuro}\else\texteuro\fi}
% Énoncés : Éducation nationale. Transcription et corrigés : APMEP, Denis Vergès et contributeurs. Découpage et classement : CoopMaths, Sébastien Lozano, Éric Elter et contributeurs. Adaptation de la présentation : Maths Explorer.
% CC BY-SA : https://coopmaths.fr/www/about/
% https://coopmaths.fr/alea/static/e3c/2020/tex/e3c_2020_00_sujet60_2.tex
\begin{document}
\noindent\textit{APMEP / CoopMaths. CC BY-SA. Présentation adaptée par Maths Explorer.}\par\medskip
\medskip

On modélise la diffusion dans le sang d'un médicament de 1 gramme par intraveineuse (fonction $f_1$, courbe représentative $C_1$) ou par voie orale (fonction $f_2$, courbe représentative $C_2$) pendant une durée de 10 heures.

Plus précisément :

\setlength\parindent{1cm}
\begin{itemize}[label=\textbullet]
\item $f_1(t)$ modélise la proportion du médicament dans le sang à l'instant $t$, où $t$ est le temps en heure après injection par intraveineuse ;
\item $f_2(t)$ modélise la proportion du médicament dans le sang à l'instant $t$, où $t$ est le temps en heure après administration par voie orale.
\end{itemize}
\setlength\parindent{0cm}

Pour tout réel $t$ de l'intervalle [0~;~10], on admet que 

\[f_1(t) = \e^{-0,57t}\quad  \text{et} \quad f_2(t) = 1,75t\e^{-t}.\]

Les courbes $C_1$ et $C_2$ de $f_1$ et $f_2$ sont représentées ci-dessous.

\begin{figure}
\begin{center}
\psset{xunit=0.7cm,yunit=7.cm,comma=true}
\begin{pspicture}(-0.5,-0.1)(10.5,1.1)
\psgrid[gridlabels=0pt,subgriddiv=1,gridwidth=0.1pt](0,0)(10,1)
\multido{\n=0.0+0.1}{11}{\psline[linewidth=0.1pt](0,\n)(10.5,\n)}
\psaxes[linewidth=1.25pt,labelFontSize=\scriptstyle,Dy=0.1]{->}(0,0)(0,0)(10.5,1.1)
\psplot[plotpoints=2000,linewidth=1.25pt,linecolor=blue]{0}{10}{2.71828 0.57 x mul neg exp}
\psplot[plotpoints=2000,linewidth=1.25pt,linecolor=red]{0}{10}{2.71828  x neg exp x mul 1.75 mul}
\uput[ur](0.2,0.9){\blue $C_1$}\uput[ur](2,0.52){\red$C_2$}
\end{pspicture}
\end{center}
\end{figure}
\smallskip

\begin{enumerate}
\item \emph{Injection par voie intraveineuse}
	\begin{enumerate}
		\item Déterminer le sens de variation de la fonction $f_1$.
		\item Résoudre graphiquement $f_1(t) < 0,1$. 
		
Interpréter la réponse dans le contexte.
	\end{enumerate}
\item \emph{Administration par voie orale}

On note $f'_2$ fonction dérivée de la fonction $f_2$.

\medskip

	\begin{enumerate}
		\item Montrer que, pour tout $t$ de [0~;~10], $f'_2(t) = 1,75(1- t)\e^{-t}$.
		\item Construire le tableau de variations de la fonction $f_2$.
		\item À quel instant $t$ la proportion de médicament dans le sang est-elle la plus élevée ?
	\end{enumerate}
\end{enumerate}

\medskip

\end{document}
